\documentclass[10pt]{article}

\usepackage[utf8]{inputenc}

\usepackage[T1]{fontenc}

\usepackage{amsmath}

\usepackage{amsfonts}

\usepackage{amssymb}

\usepackage[version=4]{mhchem}

\usepackage{stmaryrd}

\usepackage{graphicx}

\usepackage[export]{adjustbox}

\graphicspath{ {./images/} }



\begin{document}

ционных полиномов Лагранжа по узлам $2 k \pi /(2 n-1)$ (в периодич. случае) или по узлам $\cos [(2 k-1) \pi /(2 n-1-2)]$ на отрезке $[-1,1]$ (см., напр., [6]).



Вопрос о сходимости линейных положительных операторов $U_{n}^{+}$, действующих из $C[a, b]$ в $A_{n}$ или из $\tilde{C}^{\prime}$ в $T_{n}$ (в частности, операторов вица (4) с положительным ядром), решается на трех пробных функциях (см. [1]): для равномерной сходимости последовательности $U_{n}^{+}(f, t)$ к $f(t) \in C[a, b]$ или к $f(t) \in \tilde{C}$ необходимо и достаточно, чтобы это имело место для функций $1, t$, $t^{2}$ или соответственно для функций $1, \sin t, \cos t$.



Исследование погрешности приближения, доставляемой л. м. п., сводится, чаще всего, к изучению скорости сходимости $U_{N}(f, t)$ к $f(t)$, оценке погрешности через дифференциально-разностные характеристики приближаемой функции, выяснению вопроса о том, как реагирует л. м. п. на улучшение ее гладкостных свойств.



При оценке аппроксимативных свойств л. м. п. (1) естественным ориентиром служит наилучшее приближение $E\left(f, \mathfrak{l}_{N}\right)$ фу фнкции $f$ подпространством $\mathfrak{l}_{N}$. Лебега перавенство



\[

\left\|f-S_{n}(f)\right\|_{\tilde{C}} \leqslant E\left(f, T_{n}\right)_{\tilde{C}}(\ln n+3)

\]



показывает в сопоставлении с Джексона неравенством



\[

E\left(f, T_{n}\right)_{\tilde{C}} \leqslant M n-r_{\omega}\left(f^{(r)}, 1 / n\right)

\]



$(\omega(g, \delta)$ - модуль непрерывности функции $g \in \tilde{C})$, что, хотя порядок приближения суммами Фурье несколько хуже наилучшего (ln $n$ в (5) нельзя заменить константой), эти суммы реагируют на любое повышение порядка диффференцируемости приближаемой функции. Для нек-рых же л. м. п. порядок приближения не может быть выше определенной величины, сколько бы производных ни имела функция $f(t)$ (ә фф е к т н а с ы ш н и я). Так, порядок приближения линейными положительными полиномиальными операторами $U_{n}^{+}$не может быть выше $O\left(n^{-2}\right)$; для сумм Фейера порядок насыщения $O\left(n^{-1}\right)$, для сумм Бернштейна Рогозинского $O\left(n^{-2}\right)$. Интерполяционные сплайны при определенном выборе узлов склейки обеспечивают наилучший порядок погрешности приближения нө только самой функции, но и нек-рых ее первых производных в зависимости от степени многочленов, из к-рых склеен сплайн (см. [7], [8]).



В отдельных случаях для конкретных л. М. п. найдены точные или асимптотически точные оценки погрешности на классах ф্фунццй. Видимо, первый нетривиальный результат такого рода был получен А. Н. Колмогоровым, к-рый в 1935 установил, тто



\[

\sup _{f \in W^{r} K}\left\|f-S_{n}(f)\right\|_{\tilde{C}}=\frac{4 K}{\pi^{2}} \frac{\operatorname{In} n}{n^{r}}+O(n-r) \text {, }

\]



где $W^{r} K(r=1,2, \ldots)$ - класс функций $f \in \tilde{C}$, у к-рых $f^{(r-1)}(t)$ абсолютно непрерывна на $[0,2 \pi]$ и почти всюду $\left|f^{(r)}(t)\right| \leqslant K$. В дальнейшем аналогичного характера результаты были получены для сумм Фурье (и для нек-рых их средних) на других важных классах функций, задаваемых, напр., мажорантой модуля непрерывности $r$-й производной (см. [2], [6], [9]). Особый интерес представляют л. М. П. (1), в точности реализующие на том или ином классе функций верхнюю грань наилучших приближений подпространством $\mathfrak{N}_{N}$. Таким свойством для классов $W^{r} K(r=1,2, \ldots)$ обладают при определенном выборе $\lambda_{k}^{(n)}$ суммы вида (3), напр., при $r=1$ надо положить



\[

\lambda_{k}^{(n)}=\frac{k \pi}{2 n+2} \operatorname{ctg} \frac{k \pi}{2 n+2}

\]



а также интерполяционные сплайны порядка $r-1$ дефекта 1 с узлами склейки $k \pi / n(k=0,1,2, \ldots)$ (см. [4], а также Приближение функций; әкстремалыные задачи на классах фупкций, Паилучший линейный метод). Лит.: [1] К о р о в ки н ІІ. П., Линейные операторы и теория приближений, М., 1959; [2] Дं з я д ы К В. К., Введениө в теорию равномерного приближения функций полиномами, М. 1977; [3] Т и х о м и р о в B. М., Некоторые вопросы теории приближений. М.. 1976; [4] К о р н е й ч у К Н. П., Экстремальные задачи теории приближения, М., 1976; [5] $\Gamma$ о н ч ар о в В. Л., Теория интерполирования и приближения функцић, 2 изд., М., $1954 ;[6]$ Т и м а н А. Ф., Теория приближе-



\includegraphics[max width=\textwidth, center]{2023_07_08_59a3332958ad9c8210e0g-1(2)}

нов и ее приложения, пер. с англ., М., 1972; [8] С т е ч к и н С. Б., С у 6 б о т и Ю. Н., Сплайны в вычислительной математике, М., 1976; [9] С т е п а н е ц А. И., Равномерные приближения тригонометрическими полиномами. Линейные методы, K.; 1981



\includegraphics[max width=\textwidth, center]{2023_07_08_59a3332958ad9c8210e0g-1}

устанавливающие связь между дифференциально-разностными свойствами приближаемой функции и величиной (а также поведением) погрешности приближения ее тем или иным методом. Прямые теоремы (п. т.) дают оценку погрешности приближения функции $f(t)$ через ее гладкостные характеристики (наличие производных определенного порядка, модуль непрерывности или модуль гладкости самой функции $f$ или нек-рой еe производной и т. п.). В случае наилучшего приближения полиномами п. т. известны еще как теоремы Джексона [1] и их различные обобщения и уточнения (см. Джексона неравенство, Джексона теорема). Обратные теоремы (о. т.) характеризуют дифференциально-разностные свойства функций в зависимости от скорости убывания к нулю ее наилучпих (или каких-либо других) приближений. Задача получения о. т. приближения функций впервые была поставлена, а в нек-рых случаях и решена С. Н. Бернштейном [2]. Сопоставление п. т. и о. т. позволяет иногда полностью охарактеривовать класс функций, имеющих те или иные гладкостные свойства, с помощью последовательности, напр., наилучших приближений.



Наиболее проста связь между п. т. и о. т. в периодич. случае. Пусть $\tilde{c}$ - пространство непрерывных



\begin{center}

\includegraphics[max width=\textwidth]{2023_07_08_59a3332958ad9c8210e0g-1(1)}

\end{center}



\[

\|f\|_{\tilde{C}}=\max _{t}|f(t)|, \quad E\left(f, T_{n}\right)=\inf _{\varphi \in T_{n}}\|f-\varphi\|_{\tilde{C}}

\]



\begin{itemize}

  \item наилучшее приближение функции $f$ из $\tilde{C}$ подпространством $T_{n}$ триговометрич. полиномов порядка $n$, $\omega(f, \delta)$ - модуль непрерывности функции $f \in \tilde{C} ; \tilde{C}^{r}$ $(r=1,2, \ldots)$ - множество $r$ рав непрерывно диффференцируемых на всей оси функций из $\tilde{C}, \tilde{C}^{0}=\tilde{C}$. Пр я м а я те о р е м а: если $f \in \tilde{C}^{r}$, то

\end{itemize}



\[

\begin{gathered}

E\left(f, T_{n-1}\right) \leqslant \frac{M}{n^{r}} \omega\left(f(r), \frac{1}{n}\right), \\

n=1,2, \ldots, r=0,1, \ldots,

\end{gathered}

\]



где константа $M$ не зависит от $n$. Более сильное утверждение состоит в том, что можно указать последовательность линейных методов $U_{n}(n=0,1$, ...), сопоставляющих функции $f(t)$ из $\tilde{C}$ полином $U_{n}(f, t) \in T_{n}$ и таких, что для $f \in \overline{C^{r}}$ погрешность $\left\|f-U_{n}(f)\right\|_{\tilde{C}}$ оценивается правой частью (1). О б р а т н а я т е 0р е м а утверждает, что для $f \in \tilde{C}$



\[

\omega(f, \delta) \leqslant M \delta \sum_{n=1}^{[1 / \delta]} E\left(f, T_{n-1}\right), \delta>0,

\]



где $M$ - абсолютная константа, [1/8] - целая часть числа $1 / \delta$, а из сходимости при нек-ром натуральном $r$ ряда



\[

\sum_{n=1}^{\infty} n^{r-1} E\left(f, T_{n-1}\right)

\]





\end{document}